\chapter{1993年上海卷}
\section{填空题}

\begin{problem}
函数 \(y = x^2 - 2x + 3\) 的最小值是\fillinblank{}.
\end{problem}

\begin{answer}
2
\end{answer}

\begin{problem}
函数 \(y = \cos^2(\omega x)\)（\(\omega > 0\)）的最小正周期是\fillinblank{}.
\end{problem}

\begin{answer}
\(\frac{\pi}{\omega}\)
\end{answer}

\begin{problem}
设 \(z = \cos \frac{7\pi}{5} + \i \sin \frac{7\pi}{5}\)，\(\i\) 是虚数单位，则 \(z\) 的辐角主值是\fillinblank{}.
\end{problem}

\begin{answer}
\(\frac{3\pi}{5}\)
\end{answer}

\begin{problem}
\(1\text{ 名}\)老师和\(4\text{ 名}\)获奖同学排成一排照相留念，若老师不排在两端，则共有不同排法\fillinblank{}\(\text{种}\)（结果用数值表示）.
\end{problem}

\begin{answer}
72
\end{answer}

\begin{problem}
函数 \(y = \arccos x\)（\(-1 \leqslant x \leqslant 0\)）的反函数是\fillinblank{}.
\end{problem}

\begin{answer}
\(y = \cos x\)（\(\frac{\pi}{2} \leqslant x \leqslant \pi\)）
\end{answer}

\begin{problem}
\(P\) 点分 \(\overline{AB}\) 所成的比是 \(-3\)，则 \(B\) 点分 \(\overline{AP}\) 所成的比是\fillinblank{}.
\end{problem}

\begin{answer}
2
\end{answer}

\begin{problem}
已知圆台的上，下底面半径分别是 \(10 \mathrm{~cm}\) 和 \(20 \mathrm{~cm}\)，它的侧面展开后所得扇环的圆心角是 \(180^{\circ}\)，那么圆台的侧面积是\fillinblank{} \(\mathrm{cm}^{2}\)（结果中保留 \(\pi\)）.
\end{problem}

\begin{answer}
\(600\pi\)
\end{answer}

\begin{problem}
已知等比数列 \(\{a_n\}\)（\(a_n \in \R\)），\(a_1 + a_2 = 9\)，\(a_1 \cdot a_2 \cdot a_3 = 27\)，且 \(S_n = a_1 + a_2 + \cdots + a_n\)（\(n = 1\)，\(2\)，\(\cdots\)），则 \(\displaystyle\lim_{n \to \infty} S_n =\)\fillinblank{}.
\end{problem}

\begin{answer}
12
\end{answer}

\begin{problem}
在极坐标系中，过点 \(M\left(2, \frac{\pi}{2}\right)\) 且平行于极轴的直线的极坐标方程是\fillinblank{}.
\end{problem}

\begin{answer}
\(\rho \sin \theta = 2\)
\end{answer}

\begin{problem}
如图，正方体 \(ABCD - A_1B_1C_1D_1\) 中，过顶点 \(B\)，\(D\)，\(C_1\) 作截面，则二面角 \(B - DC_1 - C\) 的大小是\fillinblank{}（用反三角函数表示）.

\begin{minipage}{\linewidth}
\bitmapfigure[float=inline,max width=0.42\linewidth]{img/1993/shanghai/q10_fig1.png}
\end{minipage}
\end{problem}

\begin{answer}
\(\arccos \frac{\sqrt{3}}{3}\)
\end{answer}

\section{单选题}

\begin{problem}
函数 \(y = x + a\) 与 \(y = \log_a x\) 的图象可能是（\quad）

\choices
    {\choicebitmap[width=0.15\paperwidth]{img/1993/shanghai/q11_optA.png}}
    {\choicebitmap[width=0.15\paperwidth]{img/1993/shanghai/q11_optB.png}}
    {\choicebitmap[width=0.15\paperwidth]{img/1993/shanghai/q11_optC.png}}
    {\choicebitmap[width=0.15\paperwidth]{img/1993/shanghai/q11_optD.png}}
\end{problem}

\begin{answer}
C
\end{answer}

\begin{problem}
\((x - 1)^{9}\) 按 \(x\) 降幂排列的展开式中，系数最大的项是（\quad）

\choices
    {第4项和第5项}
    {第5项}
    {第5项和第6项}
    {第6项}
\end{problem}

\begin{answer}
B
\end{answer}

\begin{problem}
在区间 \((0, +\infty)\) 上为增函数的是（\quad）

\choices
    {\(y = \operatorname{arccot} x\)}
    {\(y = \log_{\frac{1}{2}} x\)}
    {\(y = -2^{x}\)}
    {\(y = x^{\frac{1}{3}}\)}
\end{problem}

\begin{answer}
D
\end{answer}

\begin{problem}
下列不等式中正确的是（\quad）

\choices
    {\(\arcsin \left(-\frac{1}{4}\right) < \arcsin \left(-\frac{1}{3}\right)\)}
    {\(\arccos \left(-\frac{1}{4}\right) < \arccos \frac{1}{3}\)}
    {\(\arctan \frac{1}{3} < \arcsin \frac{1}{4}\)}
    {\(\operatorname{arccot} \frac{1}{3} < \arctan \frac{1}{3}\)}
\end{problem}

\begin{answer}
C
\end{answer}

\begin{problem}
动点 \(P\) 到直线 \(x + 4 = 0\) 的距离减去它到点 \(M(2,0)\) 的距离之差等于\(2\)，则点 \(P\) 的轨迹是（\quad）

\choices
    {直线}
    {椭圆}
    {双曲线}
    {抛物线}
\end{problem}

\begin{answer}
D
\end{answer}

\begin{problem}
设有三个命题：

\begin{circlelist}
\item 底面是平行四边形的四棱柱是平行六面体；
\item 底面是矩形的平行六面体是长方体；
\item 直四棱柱是直平行六面体.
\end{circlelist}

以上命题中真命题的个数是（\quad）

\choices
    {0}
    {1}
    {2}
    {3}
\end{problem}

\begin{answer}
B
\end{answer}

\begin{problem}
在直角坐标系中平移坐标轴，把原点 \(O(0,0)\) 移到 \(O'(2,-5)\)，点 \(A\) 在新坐标系中的坐标为 \((-3,7)\)，则点 \(A\) 在原坐标系中的坐标是（\quad）

\choices
    {\((-1,2)\)}
    {\((1,-2)\)}
    {\((-5,12)\)}
    {\((5,-12)\)}
\end{problem}

\begin{answer}
A
\end{answer}

\begin{problem}
设 \(a\)，\(b\) 是两条异面直线. 在下列命题中正确的是（\quad）

\choices
    {有且仅有一条直线与 \(a\)，\(b\) 都垂直}
    {有一平面与 \(a\)，\(b\) 都垂直}
    {过直线 \(a\) 有且仅有一平面与 \(b\) 平行}
    {过空间中任一点必可作一条直线与 \(a\)，\(b\) 都相交}
\end{problem}

\begin{answer}
C
\end{answer}

\begin{problem}
函数 \(y = \cos \left(2x + \frac{\pi}{2}\right) \) 的图象的一条对称轴方程是（\quad）

\choices
    {\(x = -\frac{\pi}{2}\)}
    {\(x = -\frac{\pi}{4}\)}
    {\(x = \frac{\pi}{8}\)}
    {\(x = \pi\)}
\end{problem}

\begin{answer}
B
\end{answer}

\begin{problem}
\(a + b > 2c^{2}\) 的一个充分条件是（\quad）

\choices
    {\(a > c\) 或 \(b > c\)}
    {\(a > c\) 且 \(b < c\)}
    {\(a > c\) 且 \(b > c\)}
    {\(a > c\) 或 \(b < c\)}
\end{problem}

\begin{answer}
C
\end{answer}

\section{解答题}

\begin{problem}
已知角 \(\alpha\) 的顶点与直角坐标系的原点重合，始边在 \(x\) 轴的正半轴上，终边经过点 \(P(-1,2)\)，求 \(\sin \left(2\alpha + \frac{2\pi}{3}\right)\) 的值.
\end{problem}

\begin{solution}
【法一】 因为 \(r = |OP| = \sqrt{(-1)^2 + 2^2} = \sqrt{5}\)，所以 \(\sin \alpha = \frac{2}{\sqrt{5}}\)，\(\cos \alpha = -\frac{1}{\sqrt{5}}\). 因为 \(\sin 2\alpha = 2\sin \alpha \cos \alpha = -\frac{4}{5}\)，\(\cos 2\alpha = 1 - 2\sin^2 \alpha = -\frac{3}{5}\). 所以 \(\sin \left(2\alpha + \frac{2\pi}{3}\right) = \sin 2\alpha \cdot \cos \frac{2\pi}{3} + \cos 2\alpha \cdot \sin \frac{2\pi}{3} = \frac{4 - 3\sqrt{3}}{10}\).

【法二】 因为 \(\tan \alpha = -2\)，所以 \(\sin 2\alpha = \frac{2\tan \alpha}{1 + \tan^2 \alpha} = -\frac{4}{5}\)，\(\cos 2\alpha = \frac{1 - \tan^2 \alpha}{1 + \tan^2 \alpha} = -\frac{3}{5}\). 以下同解一.
\end{solution}

\begin{problem}
已知 \(z = \frac{a - \i}{1 - \i}\)，其中 \(\i\) 为虚数单位，\(a > 0\)，复数 \(w = z(z + \i)\) 的虚部减去它的实部所得的差等于 \(\frac{3}{2}\)，求复数 \(w\) 的模.
\end{problem}

\begin{solution}
\(w = \frac{a - \i}{1 - \i}\left(\frac{a - \i}{1 - \i} + \i\right) = \frac{(a - \i)(a + 1)}{(1 - \i)^2} = \frac{a + 1}{2}(1 + a\i)\). 由已知条件，\(\frac{a + 1}{2}a - \frac{a + 1}{2} = \frac{3}{2}\)，即 \(a^2 = 4\). 由已知 \(a > 0\)，故 \(a = 2\). 于是 \(|w| = \left|\frac{3}{2} + 3\i\right| = \frac{3}{2}\sqrt{5}\).
\end{solution}

\begin{problem}
抛物线 \(y = -\frac{x^2}{2}\) 与过点 \(M(0,-1)\) 的直线 \(l\) 相交于 \(A\)，\(B\) 两点，\(O\) 为坐标原点. 若直线 \(OA\) 与 \(OB\) 的斜率之和为 \(1\)，求直线 \(l\) 的方程.
\end{problem}

\begin{solution}
【法一】 设点 \(A\) 和点 \(B\) 的坐标分别为 \((x_1,y_1)\) 和 \((x_2,y_2)\). 由条件可设直线 \(l\) 的方程为 \(y = kx - 1\). 设直线 \(OA\) 和 \(OB\) 的斜率分别为 \(k_{OA}\) 和 \(k_{OB}\)，则由条件有
\(k_{OA} + k_{OB} = \frac{y_1}{x_1} + \frac{y_2}{x_2} = 1\).
由于 \(y_1 = -\frac{x_1^2}{2}\)，\(y_2 = -\frac{x_2^2}{2}\)，故代入上式后得 \(-\frac{1}{2}(x_1 + x_2) = 1\). 于是
\[
k = \frac{y_2 - y_1}{x_2 - x_1} = -\frac{\left(-\frac{x_2^2}{2}\right) - \left(-\frac{x_1^2}{2}\right)}{x_2 - x_1} = -\frac{1}{2}(x_1 + x_2) = 1
\]
因此直线 \(l\) 的方程为 \(y = x - 1\).

【法二】由法一中的条件式，直线方程代入抛物线方程得 \(x^2 + 2kx - 2 = 0\).
方程有实数解 \(x_1\) 及 \(x_2\)，由韦达定理，有 \(x_1 + x_2 = -2k\)，\(x_1x_2 = -2\).
由于 \(y_1 = kx_1 - 1\)，\(y_2 = kx_2 - 1\)，于是条件式成为
\[
1 = \frac{kx_1 - 1}{x_1} + \frac{kx_2 - 1}{x_2} = 2k - \frac{x_1 + x_2}{x_1x_2} = k
\]
因此直线 \(l\) 的方程为 \(y = x - 1\).

【法三】由法二中的方程，方程的解为 \(x_1 = -k - \sqrt{k^2 + 2}\)，\(x_2 = -k + \sqrt{k^2 + 2}\).
由于 \(y_1 = -\frac{x_1^2}{2}\)，\(y_2 = -\frac{x_2^2}{2}\)，于是条件式成为
\[
1 = -\frac{x_1}{2} - \frac{x_2}{2} = -\frac{1}{2}\left(-k - \sqrt{k^2 + 2}\right) - \frac{1}{2}\left(-k + \sqrt{k^2 + 2}\right) = k
\]
因此直线 \(l\) 的方程为 \(y = x - 1\).
\end{solution}

\begin{problem}
如图，已知二面角 \(\alpha - PQ - \beta\) 为 \(60^{\circ}\)，点 \(A\) 和点 \(B\) 分别在平面 \(\alpha\) 和平面 \(\beta\) 上，点 \(C\) 在棱 \(PQ\) 上，\(\angle ACP = \angle BCP = 30^{\circ}\)，\(CA = CB = a\).

\begin{enumerate}
\item 求证 \(AB \perp PQ\).
\item 求点 \(B\) 到平面 \(\alpha\) 的距离.
\item 设 \(R\) 是线段 \(CA\) 上的一点，直线 \(BR\) 与平面 \(\alpha\) 所成角的大小为 \(45^{\circ}\)，求线段 \(CR\) 的长度.
\end{enumerate}

\bitmapfigure[float=inline,max width=0.48\linewidth]{img/1993/shanghai/q24_fig1.png}
\end{problem}

\begin{solution}
【法一】
\begin{enumerate}
\item 在平面 \(\beta\) 内作 \(BD \perp PQ\)，垂足为 \(D\)，连接 \(AD\). 因 \(\angle BCD = \angle ACD\)，且 \(BC = AC\)，\(DC\) 为公共边，故 \(\triangle BCD \cong \triangle ACD\)，所以 \(AD \perp PQ\). 于是 \(PQ \perp\) 平面 \(ABD\)，故 \(AB \perp PQ\).

\item 因 \(BD \perp PQ\)，\(AD \perp PQ\)，故 \(\angle ADB\) 为二面角 \(\alpha - PQ - \beta\) 的平面角，\(\angle ADB = 60^{\circ}\). 在 \(\mathrm{Rt}\triangle BCD\) 中，因为 \(BC = a\)，\(\angle BCD = 30^{\circ}\)，所以 \(BD = \frac{1}{2}a\). 过 \(B\) 作 \(BH \perp AD\)，\(H\) 为垂足. 因为 \(PQ \perp\) 平面 \(ABD\)，所以平面 \(ABD \perp \alpha\)，\(BH \perp \alpha\). 于是点 \(B\) 到平面 \(\alpha\) 的距离为 \(BH = BD \cdot \sin 60^{\circ} = \frac{\sqrt{3}}{4}a\).

\item 连接 \(HR\). 因 \(BH \perp \alpha\)，故 \(\angle BRH\) 是 \(BR\) 与平面 \(\alpha\) 所成的角，\(\angle BRH = 45^{\circ}\). 因 \(BH = \frac{\sqrt{3}}{4}a\)，故 \(BR = \frac{\sqrt{6}}{4}a\). 在 \(\triangle ABD\) 中，\(AD = BD\)，\(\angle ADB = 60^{\circ}\)，所以 \(\triangle ABD\) 为正三角形，\(AB = BD = \frac{a}{2}\). 在 \(\triangle ABC\) 中，\(AC = BC = a\)，由余弦定理得 \(\cos \angle BCA = \frac{a^2 + a^2 - \left(\frac{a}{2}\right)^2}{2 \cdot a \cdot a} = \frac{7}{8}\).
在 \(\triangle BCR\) 中，设 \(x = CR\)，由余弦定理得 \(\left(\frac{\sqrt{6}}{4}a\right)^2 = x^2 + a^2 - 2ax \cdot \frac{7}{8}\)，
即 \(8x^2 - 14ax + 5a^2 = 0\). 所以 \(x_1 = \frac{1}{2}a\)，\(x_2 = \frac{5}{4}a\).
由于 \(\frac{5}{4}a > a\)，故 \(x_2 = \frac{5}{4}a\) 舍去. 因此 \(CR\) 的长为 \(\frac{1}{2}a\).

\end{enumerate}

【法二】
\begin{enumerate}
\item 同法一.

\item 同法一.

\item 连接 \(HR\). 因 \(BH \perp \alpha\)，故 \(\angle BRH\) 是 \(BR\) 与平面 \(\alpha\) 所成的角，\(\angle BRH = 45^{\circ}\). \(HR = BH = \frac{\sqrt{3}}{4}a\). 因 \(\triangle ADC\) 为 \(\mathrm{Rt}\triangle\)，\(\angle DAC = 60^{\circ}\). 在正三角形 \(ABD\) 中，\(BH \perp AD\)，所以 \(AH = \frac{1}{2}AD = \frac{1}{4}a\). 在 \(\triangle ARH\) 中，设 \(t = AR\)，由余弦定理得 \(\left(\frac{\sqrt{3}}{4}a\right)^2 = \left(\frac{1}{4}a\right)^2 + t^2 - 2 \cdot \frac{1}{4}a \cdot t \cdot \frac{1}{2}\)，
即 \(8t^2 - 2at - a^2 = 0\). 所以 \(t_1 = \frac{1}{2}a\)，\(t_2 = -\frac{1}{4}a\).
因为 \(-\frac{1}{4}a < 0\)，故 \(t_2 = -\frac{1}{4}a\) 舍去. 因此 \(CR = a - \frac{1}{2}a = \frac{1}{2}a\).
\end{enumerate}
\end{solution}

\begin{problem}
设数列 \(\{a_n\}\) 的前 \(n\) 项和 \(S_n = na + n(n - 1)b\)（\(n = 1\)，\(2\)，\(\cdots\)），\(a\)，\(b\) 是常数且 \(b \neq 0\).

\begin{enumerate}
\item 证明 \(\{a_n\}\) 是等差数列.
\item 证明以 \(\left(a_n, \frac{S_n}{n} - 1\right)\) 为坐标的点 \(P_n\)（\(n = 1\)，\(2\)，\(\cdots\)）都落在同一条直线上，并写出此直线的方程.
\item 设 \(a = 1\)，\(b = \frac{1}{2}\)，\(C\) 是以 \((r,r)\) 为圆心，\(r\) 为半径的圆（\(r > 0\)）. 求使得点 \(P_1\)，\(P_2\)，\(P_3\) 都落在圆 \(C\) 外时，\(r\) 的取值范围.
\end{enumerate}
\end{problem}

\begin{solution}
\begin{enumerate}
\item 由条件得 \(a_1 = S_1 = a\). 当 \(n \geqslant 2\) 时，有
\[
a_n = S_n - S_{n - 1} = \left[na + n(n - 1)b\right] - \left[(n - 1)a + (n - 1)(n - 2)b\right] = a + 2(n - 1)b
\]
因此，当 \(n \geqslant 2\) 时，\(a_n - a_{n - 1} = \left[a + 2(n - 1)b\right] - \left[a + 2(n - 2)b\right] = 2b\).
因此，\(\{a_n\}\) 是以 \(a\) 为首项，\(2b\) 为公差的等差数列.

\item 因为 \(b \neq 0\)，对于 \(n \geqslant 2\)，有
\[
\frac{\left(\frac{S_n}{n} - 1\right) - \left(\frac{S_1}{1} - 1\right)}{a_n - a_1}
= \frac{\frac{na + n(n - 1)b}{n} - a}{a + 2(n - 1)b - a}
= \frac{(n - 1)b}{2(n - 1)b}
= \frac{1}{2}
\]
因此，所有的点 \(P_n\left(a_n,\frac{S_n}{n} - 1\right)\)（\(n = 1\)，\(2\)，\(\cdots\)）都落在通过 \(P_1(a,a - 1)\) 且以 \(\frac{1}{2}\) 为斜率的直线上. 此直线方程为 \(y - (a - 1) = \frac{1}{2}(x - a)\)，即 \(x - 2y + a - 2 = 0\).

\item 当 \(a = 1\)，\(b = \frac{1}{2}\) 时，\(P_n\) 的坐标为 \(\left(n,\frac{n - 1}{2}\right)\). 使 \(P_1(1,0)\)，\(P_2\left(2,\frac{1}{2}\right)\)，\(P_3(3,1)\) 都落在圆 \(C\) 外的条件是
\[
\begin{cases}
(r - 1)^2 + r^2 > r^2,\\
(r - 2)^2 + \left(r - \frac{1}{2}\right)^2 > r^2,\\
(r - 3)^2 + (r - 1)^2 > r^2.
\end{cases}
\]
即等价于
\[
\begin{cases}
(r - 1)^2 > 0,\\
r^2 - 5r + \frac{17}{4} > 0,\\
r^2 - 8r + 10 > 0.
\end{cases}
\]
由第一个不等式，得 \(r \neq 1\)；由第二个不等式，得 \(r < \frac{5}{2} - \sqrt{2}\) 或 \(r > \frac{5}{2} + \sqrt{2}\)；由第三个不等式，得 \(r < 4 - \sqrt{6}\) 或 \(r > 4 + \sqrt{6}\). 再注意到 \(r > 0\)，且
\[
1 < \frac{5}{2} - \sqrt{2} < 4 - \sqrt{6} < \frac{5}{2} + \sqrt{2} < 4 + \sqrt{6}
\]
故使 \(P_1\)，\(P_2\)，\(P_3\) 都落在圆 \(C\) 外时，\(r\) 的取值范围是
\((0,1) \cup \left(1,\frac{5}{2} - \sqrt{2}\right) \cup (4 + \sqrt{6}, +\infty)\).
\end{enumerate}
\end{solution}

\begin{problem}
如图，\(P\) 为椭圆 \(\frac{x^2}{a^2} + y^2 = 1\) 上的一个动点，它与长轴端点不重合，\(a \geqslant \sqrt{2}\)，点 \(F_1\) 和点 \(F_2\) 分别是双曲线 \(\frac{x^2}{a^2} - y^2 = 1\) 的左焦点和右焦点，\(\varphi = \angle F_1PF_2\).

\begin{enumerate}
\item 求 \(\tan\varphi\) 的表达式（用 \(a\) 及描述点 \(P\) 位置的一个变量来表示）.
\item 当 \(a\) 固定时，求 \(\varphi\) 的最小值 \(\varphi_0\).
\item 当 \(a\) 在区间 \([\sqrt{2}, \sqrt{3}]\) 上变化时，求 \(\varphi_0\) 的取值范围.
\end{enumerate}

\bitmapfigure[float=inline,max width=0.46\linewidth]{img/1993/shanghai/q26_fig1.png}
\end{problem}

\begin{solution}
【法一】
\begin{enumerate}
\item 双曲线的半焦距 \(c = \sqrt{a^2 + 1}\)，焦点为 \(F_1(-\sqrt{a^2 + 1},0)\) 和 \(F_2(\sqrt{a^2 + 1},0)\). 设 \(P(x_0,y_0)\) 为椭圆上一点，与长轴的端点 \(A_1\) 及 \(A_2\) 不重合，由对称性，不妨设 \(0 < y_0 \leqslant 1\).

作 \(PH \perp F_1F_2\)，垂足为 \(H\). 记 \(\alpha = \angle F_1PH\)，\(\beta = \angle HPF_2\)，则 \(\varphi = \alpha + \beta\). 因 \(\tan \alpha = \frac{x_0 + \sqrt{a^2 + 1}}{y_0}\)，\(\tan \beta = \frac{\sqrt{a^2 + 1} - x_0}{y_0}\)，
故
\[
\tan \varphi
= \tan(\alpha + \beta)
= \frac{\tan \alpha + \tan \beta}{1 - \tan \alpha \cdot \tan \beta}
= \frac{2\sqrt{a^2 + 1} \cdot y_0}{x_0^2 + y_0^2 - (a^2 + 1)}
\]
由于 \(P(x_0,y_0)\) 在椭圆上，\(x_0^2 = a^2(1 - y_0^2)\)，故 \(\tan \varphi = -\frac{2\sqrt{a^2 + 1} \cdot y_0}{(a^2 - 1)y_0^2 + 1}\).

\item 由于 \(a \geqslant \sqrt{2}\)，故 \(a^2 - 1 \geqslant 1\)，\(\tan \varphi < 0\)，\(\frac{\pi}{2} < \varphi < \pi\). 又
\[
\tan(\pi - \varphi)
= \frac{2\sqrt{a^2 + 1} \cdot y_0}{(a^2 - 1)y_0^2 + 1}
= \frac{2\sqrt{a^2 + 1}}{(a^2 - 1)y_0 + \frac{1}{y_0}}
\leqslant \sqrt{\frac{a^2 + 1}{a^2 - 1}}
\]
当且仅当 \((a^2 - 1)y_0 = \frac{1}{y_0}\)，即 \(y_0 = \frac{1}{\sqrt{a^2 - 1}} \leqslant 1\) 时取“\(=\)”号. 由于 \(0 < \pi - \varphi < \frac{\pi}{2}\)，且正切函数在 \(\left(0,\frac{\pi}{2}\right)\) 内是增函数，故有 \(\pi - \varphi \leqslant \arctan\sqrt{\frac{a^2 + 1}{a^2 - 1}}\).
于是 \(\varphi \geqslant \pi - \arctan\sqrt{\frac{a^2 + 1}{a^2 - 1}}\)，
从而 \(\varphi_0 = \pi - \arctan\sqrt{\frac{a^2 + 1}{a^2 - 1}}\).

\item \(\varphi_0 = \pi - \arctan\sqrt{1 + \frac{2}{a^2 - 1}}\). 由 \(\sqrt{2} \leqslant a \leqslant \sqrt{3}\)，故 \(2 \leqslant a^2 \leqslant 3\)，\(1 \leqslant a^2 - 1 \leqslant 2\)，于是
\(1 \leqslant \frac{2}{a^2 - 1} \leqslant 2\)，\(2 \leqslant 1 + \frac{2}{a^2 - 1} \leqslant 3\)，
\(\sqrt{2} \leqslant \sqrt{1 + \frac{2}{a^2 - 1}} \leqslant \sqrt{3}\).
由反正切函数单调性，
\[
\arctan\sqrt{2} \leqslant \arctan\sqrt{1 + \frac{2}{a^2 - 1}} \leqslant \arctan\sqrt{3} = \frac{\pi}{3}
\]
于是 \(\frac{2\pi}{3} \leqslant \varphi_0 \leqslant \pi - \arctan\sqrt{2}\).
故 \(\varphi_0\) 的取值范围是 \(\left[\frac{2\pi}{3}, \pi - \arctan\sqrt{2}\right]\).

\end{enumerate}

【法二】
\begin{enumerate}
\item 双曲线的半焦距 \(c = \sqrt{a^2 + 1}\)，焦点为 \(F_1(-\sqrt{a^2 + 1},0)\) 和 \(F_2(\sqrt{a^2 + 1},0)\). 设 \(P(a\cos \theta,\sin \theta)\) 为椭圆上一点，与长轴的端点 \(A_1\) 及 \(A_2\) 不重合，由对称性，不妨设 \(0 < \theta \leqslant \frac{\pi}{2}\). 记 \(\alpha = \angle PF_2x\)，\(\beta = \angle PF_1x\)，故 \(\varphi = \alpha - \beta\). 因 \(PF_1\) 的斜率 \(k_1 = \frac{\sin \theta}{a\cos \theta + \sqrt{a^2 + 1}}\)，
\(PF_2\) 的斜率 \(k_2 = \frac{\sin \theta}{a\cos \theta - \sqrt{a^2 + 1}}\)，
故 \(\tan \varphi = \tan(\alpha - \beta) = -\frac{2\sqrt{a^2 + 1} \cdot \sin \theta}{(a^2 - 1)\sin^2 \theta + 1}\).

\item 由于 \(a \geqslant \sqrt{2}\)，故 \(a^2 - 1 \geqslant 1\)，\(\tan \varphi < 0\)，\(\frac{\pi}{2} < \varphi < \pi\). 故
\[
\tan(\pi - \varphi) = \frac{2\sqrt{a^2 + 1} \cdot \sin \theta}{(a^2 - 1)\sin^2 \theta + 1}
= \frac{2\sqrt{a^2 + 1}}{(a^2 - 1)\sin \theta + \frac{1}{\sin \theta}}
\leqslant \sqrt{\frac{a^2 + 1}{a^2 - 1}}
\]
当且仅当 \((a^2 - 1)\sin \theta = \frac{1}{\sin \theta}\)，即 \(\sin \theta = \frac{1}{\sqrt{a^2 - 1}} \leqslant 1\) 时取“\(=\)”号.

\item 同法一.
\end{enumerate}
\end{solution}
